\section{Experimental design}
\label{sec:design}

\subsection{Factors}

\begin{itemize}
  \item \textbf{Architecture}: FEDOT.LLM, AutoDS-Tools, Terminus-2.
  \item \textbf{\Ktool{}}: on/off, crossed on AutoDS-Tools and Terminus-2
        (Sections~\ref{sec:grid} and~\ref{sec:kdiscterm}). Terminus-2 has no
        prescription of its own to switch, so the text is appended to the task
        file instead; provisioning, the nearest thing to a library manipulation
        that does not touch the instruction, is varied separately in
        Section~\ref{sec:matched}.
  \item \textbf{\Kdisc{}}: on/off, crossed on AutoDS-Tools and
        Terminus-2 (Sections~\ref{sec:grid} and~\ref{sec:kdiscterm}).
  \item \textbf{Backbone model}: \texttt{gemma-4-31b-it} throughout, with
        \texttt{gemma-4-26b-a4b-it} below it and \texttt{glm-4.7} above it,
        crossed on AutoDS-Tools and Terminus-2 (Section~\ref{sec:modelaxis}). The
        upper point leaves the family because that family has nothing larger, so
        this axis varies backbone capability and not parameter count.
  \item \textbf{Task protocol}: greenfield (nothing supplied) versus
        improve-a-baseline (a working training script supplied).
\end{itemize}

\subsection{Benchmarks}

\paragraph{TabReD} Eight real-world tabular tasks with temporal
splits~\citep{rubachev2024tabred}. The benchmark ships no agentic protocol, so
tasks are solved from scratch: this is our greenfield condition. Its value here
is the published reference table of eighteen methods, spanning gradient boosting
implementations, tabular deep learning architectures, ensembles and
retrieval-augmented models, tuned with Optuna on the official validation split
and averaged over fifteen seeds. Metrics are ROC-AUC for classification and RMSE
in raw target units for regression.

\paragraph{MLAgentBench} Ten tasks spanning tabular, text, vision and graph
data~\citep{huang2024mlagentbench}. Each supplies a working but weak
\texttt{train.py} and asks the agent to improve it: this is our
improve-a-baseline condition. FEDOT.LLM participates only in the two tabular
tasks, for the structural reason given in Section~\ref{sec:scope}.

This benchmark carries no reference table. Unlike TabReD it ships no tuned
external methods, and none exist in the literature at the granularity we would
need, so there is no range against which a score can be normalised and no
external ceiling to fall short of. Every quantity we report on MLAgentBench is
therefore a within-study contrast: one branch of a system against another, one
architecture against another, or one edition of an instruction layer against
another, always under conditions we held fixed ourselves. The normalised
$0$-to-$1$ language of Section~\ref{sec:ceiling} belongs to TabReD alone and
should not be carried across. What MLAgentBench contributes is the second
starting point---a working baseline to improve rather than an empty
directory---and the signs and failures that condition produces.

\paragraph{Scientific case studies} Three published studies with author-reported
baselines, used in Section~\ref{sec:cases} to test external validity outside
benchmark conditions. Their instructions carry the heaviest prescription of any
condition in this paper: library, featurisation, leakage handling and the
baseline to beat. They are read as the maximal-prescription end of the axis
rather than as an unaided-agent condition.

\subsection{Common execution environment}

All runs execute inside the Harbor framework~\citep{harbor2026}, which records for every trial the
submission, the scorer output, the number of attempts and errors, wall-clock
time, token counts and monetary cost. Reporting cost and attempts alongside
accuracy lets us treat efficiency as a second axis rather than an afterthought.

\subsection{Provisioning}

Because a tool prescription presupposes the tool (Section~\ref{sec:axis}), we
state the container contents per arm rather than assuming a neutral environment.
The Terminus-2 and FEDOT.LLM runs reported here used the benchmark's stock image:
\texttt{numpy}, \texttt{pandas} and \texttt{scikit-learn} on Python~3.13. The
AutoDS-Tools arm also requires the libraries its prescription names, on
Python~3.12, which is the environment its earlier runs also had. Data, splits,
metrics and the backbone model are identical across arms; the installed library
set is not, and we report it as a factor.

The same honesty is owed about the instruction. Each system enters the
leaderboard in the configuration it is shipped and used in, which is
not the same configuration for all three: Terminus-2 receives the task statement
alone, FEDOT.LLM receives it through a pipeline whose tool choice is welded in,
and AutoDS-Tools receives it with its prescribed-knowledge layer applied, that
layer being the system's default. A row of Table~\ref{tab:leaderboard} therefore
compares deployed systems rather than isolated architectures. Two further arms
separate the strands, AutoDS-Tools with the layer disabled on the common
adapter and Terminus-2 on the enriched image, and each is a single sweep of
the same eight tasks. We ran both, and report them in
Sections~\ref{sec:grid} and~\ref{sec:matched}. The first arm was run as a full
$2\times2$ rather than as a single off-switch, since the layer has two halves
that can be disabled independently.

\subsection{Noise floor}

Rather than assume a significance threshold, we measure one. Terminus-2 and
AutoDS-Tools were each run three times on every TabReD task;
The two systems do not share a floor.
Terminus-2's relative standard deviation ranges from $0.03\%$ to $0.94\%$;
AutoDS-Tools reproduces itself an order of magnitude more tightly, from $0.00\%$
to $0.07\%$. On six of the eight tasks for Terminus-2 and seven of the eight for
AutoDS-Tools, two attempts are bit-identical, meaning the agent reproduced the
same script, so the repeats carry less independent information than their count
suggests.

We therefore treat relative differences above $1\%$ as real when comparing the
two systems, which takes the threshold from the noisier of the pair and is the
conservative choice there. That threshold would be far too coarse for comparing
configurations of one system against each other: for AutoDS-Tools the
appropriate figure is set by its worst task, \texttt{maps-routing} at
$0.07\%$, while four of its eight tasks reproduce to every digit across all
three attempts and a fifth does so across the two it has. We use $0.1\%$ for those comparisons. One caveat applies to both figures.
Attempts that agree to every printed digit, of which there are several, are not
independent samples of the agent: on these tasks the submission is a
deterministic library call, so any two attempts that reach the same call return
the same number (Section~\ref{sec:kdiscterm}). The floors below therefore
measure library determinism as much as agent stability, which makes them
optimistic in exactly the way a noise floor should not be. The MLAgentBench cells are no longer single runs. We repeated them at three
attempts per cell---sixty trials across the two branches of the layer ablation,
twenty-four more for the second architecture, and twenty-four more for the second
edition of the layer (Sections~\ref{sec:mechanism}, \ref{sec:scope} and
\ref{sec:threats}). Two tasks resisted repetition on hardware without an
accelerator and keep their single-run figures (Section~\ref{sec:scope}).

The repeats supply the noise floor this benchmark previously lacked, and it is
nothing like the tabular one. Where the TabReD floor is a fraction of a percent,
three attempts of one configuration on \texttt{ogbn-arxiv} span $0.2237$ to
$0.4790$ and on \texttt{fathomnet} $0.0844$ to $0.3846$; the second architecture
returned $0.1144$, $0.3303$ and $0.5569$ on \texttt{fathomnet} in an earlier job
and $0.0844$, $0.2629$ and $0.3846$ in a later one. A difference on this
benchmark is worth reading only when it exceeds a spread of that size, which is
why Section~\ref{sec:mechanism} rests on cells that go to zero rather than on
cells that move a little.

Repetition also changed how we count a lost trial, and the change matters beyond
this paper. On TabReD a trial that reaches the task ceiling has written nothing,
so excluding it loses no information. On MLAgentBench that inference is false:
the verifier scores whatever \texttt{submission.csv} exists at the moment the
agent is stopped, and trials that exhausted their budget frequently submitted
usable work---\texttt{clrs} scored $0.4118$ and $0.4404$ on two such attempts and
\texttt{feedback} $0.6492$ on a third. Exhausting the budget and returning
nothing are separate events here, and we count them separately throughout.

\subsection{What is not crossed}

Two limitations of the design are deliberate. FEDOT.LLM cannot enter the
non-tabular cells at all. And the AutoDS-Tools TabReD runs reported in earlier
drafts were computed on a different data slice; they are excluded from
Table~\ref{tab:leaderboard}, which uses a re-run on the common adapter. Both
sets are released, because the difference between them is itself the measured
price of the data preparation.
